\sequence{Algebra, equations and inequalities}{Algebra}{4--20 January 2027}
\section{Brackets and distributivity}
\begin{methode}[Removing brackets]
After a $+$ sign, removing the brackets keeps the signs of the terms. After a $-$ sign, every term inside changes sign. This rule applies to a sum or a difference of terms.
\end{methode}
\ex{Remove the brackets and simplify: $A=x+(2+3x)$; $B=x-(2+3x)$; $C=-2x+(-4+5x)$; $D=-2x-(-4+5x)$.}
\reponse[3.2]{$A=4x+2$; $B=-2x-2$; $C=3x-4$; $D=-7x+4$.}
\begin{cours}[Distributivity]
For all real numbers $a,b,c,d$:
\[b(c+d)=bc+bd,\qquad(a+b)(c+d)=ac+ad+bc+bd.\]
\end{cours}
\ex{Expand and simplify:
\begin{align*}
A&=-4(8x+7)+(-8x+5)(5x+9),\\
B&=(6x+9)(-x+7)+(3x+5)(-9x+6),\\
C&=-9(-x-7)-(-x+3)(-8x-8),\\
D&=-5(-6x-7)+(x-7)(5x+8).
\end{align*}}
\reponse[5]{Expand each product, then collect terms of the same degree:
\begin{align*}
A&=-32x-28-40x^2-72x+25x+45=-40x^2-79x+17,\\
B&=-6x^2+33x+63-27x^2-27x+30=-33x^2+6x+93,\\
C&=9x+63-(8x^2-16x-24)=-8x^2+25x+87,\\
D&=30x+35+5x^2-27x-56=5x^2+3x-21.
\end{align*}}

\fiche{Factorising with a common factor}
\begin{cours}[Property]
$ka+kb=k(a+b)$ and $ka-kb=k(a-b)$. Factorising turns a sum into a product; expanding does the reverse.
\end{cours}
\ex{Factorise:
\begin{align*}
A&=15x^2-21x,\\B&=3(2+3x)-(5+2x)(2+3x),\\
C&=(4x-1)(x+6)+(4x-1),\\D&=(1-6x)^2-(1-6x)(2+5x).
\end{align*}}
\reponse[6]{
$A=3x(5x-7)$.\\
$B=(2+3x)[3-(5+2x)]=(2+3x)(-2-2x)=-2(2+3x)(1+x)$.\\
$C=(4x-1)[(x+6)+1]=(4x-1)(x+7)$.\\
$D=(1-6x)[(1-6x)-(2+5x)]=(1-6x)(-1-11x)=(6x-1)(11x+1)$.
Check the factors by expanding again.}
\ex{Choose a useful form of $f(x)=(x-2)(x+3)$ to calculate $f(2)$, $f(0)$ and find its zeros.}
\reponse[3.2]{The factorised form immediately gives $f(2)=0$ and the zeros $2,-3$. At $0$, $f(0)=(-2)\times3=-6$; the expanded form $x^2+x-6$ also gives this result. The choice depends on the question.}

\fiche{Special identities: expanding}
\begin{cours}[Identities]
\[(a+b)^2=a^2+2ab+b^2,\qquad(a-b)^2=a^2-2ab+b^2,\]
\[(a+b)(a-b)=a^2-b^2.\]
\end{cours}
\preuve[3.2]{Expand $(a+b)(a+b)$: $a^2+ab+ba+b^2=a^2+2ab+b^2$. Replacing $b$ by $-b$ gives the second identity. Finally $(a+b)(a-b)=a^2-ab+ab-b^2=a^2-b^2$.}
\ex{Expand $(x+3)^2$, $(3x-4)^2$, $(x-3)(x+3)$, $(2x+3)(2x-3)$, $(4-3x)(3x+4)$.}
\reponse[3.2]{$x^2+6x+9$; $9x^2-24x+16$; $x^2-9$; $4x^2-9$; $16-9x^2$. In the last product, reorder $3x+4$ as $4+3x$.}
\ex{Expand and simplify $A=(2x-3)^2+(x+5)(3-x)$ and $B=2(x+3)+(2x+3)(2x-3)$.}
\reponse[4]{$A=4x^2-12x+9-x^2-2x+15=3x^2-14x+24$. $B=2x+6+4x^2-9=4x^2+2x-3$.}

\fiche{Special identities: factorising}
\begin{methode}
Spot the two squares; for three-term forms, check the double product $2ab$. For a difference of two squares, use $a^2-b^2=(a-b)(a+b)$.
\end{methode}
\ex{Factorise:
\[A=x^2-2x+1,\quad B=25+16x^2-40x,\quad C=4x^2+12x+9,\]
\[D=1-49x^2,\quad E=9x^2-4,\quad F=(2x+3)^2-64.\]}
\reponse[6]{
$A=(x-1)^2$. $B=(4x-5)^2$. $C=(2x+3)^2$.\\
$D=(1-7x)(1+7x)$. $E=(3x-2)(3x+2)$.\\
$F=[(2x+3)-8][(2x+3)+8]=(2x-5)(2x+11)$.
For the first three, check the middle term; for $F$, treat $2x+3$ as a single block.}
\ex[Checking an identity]{Are the equalities $(a+b)^2=a^2+b^2$ and $a^2+b^2=(a-b)(a+b)$ true for all real numbers?}
\reponse[3.2]{No: for $a=b=1$, the first would give $4=2$ and the second $2=0$. The missing term in the first is $2ab$; the product in the second equals $a^2-b^2$.}
\rituel{Calculate $101^2$ mentally; give a direction vector of $2x+3y=1$; determine $[-1;3]\cap[2;5]$.}{$10201$; $(-3;2)$; $[2;3]$.}

\fiche{Integer powers}
\begin{cours}[Definition and rules]
For $n\geq1$, $a^n$ is the product of $n$ factors equal to $a$. For $a\ne0$, $a^0=1$ and $a^{-n}=1/a^n$. In this course, $0^0$ is not defined.
For $a,b\ne0$ and integers $m,n$:
\[a^ma^n=a^{m+n},\quad\frac{a^m}{a^n}=a^{m-n},\quad(a^m)^n=a^{mn},\]
\[(ab)^n=a^nb^n,\qquad\left(\frac ab\right)^n=\frac{a^n}{b^n}.\]
For strictly positive exponents, some bases may be zero if the expressions remain defined.
\end{cours}
\ex{Simplify $A=2^4\times2^3$, $B=(-5)^{12}/(-5)^3$, $C=(7^3)^6\times7^{-5}$, $D=7^{24}\times3^2\times2^3$. Also calculate $(-3)^5$, $5^1$ and $99999^0$.}
\reponse[4]{$A=2^7=128$. $B=(-5)^9=-1953125$. $C=7^{18-5}=7^{13}$. $D=72\times7^{24}$: exponents of different bases cannot be added. The last three values are $-243$, $5$, $1$.}
\vigilance{The definition as a repeated product does not hold for a negative exponent. $-3^2=-9$, whereas $(-3)^2=9$.}
\ex{Solve $2^n>100$ mentally for the smallest natural number $n$. Write a program that finds it.}
\reponse[4]{$2^6=64\leq100<128=2^7$, so $n=7$. Initialise $n=0$, $p=1$; while $p\leq100$, replace $p$ by $2p$ and $n$ by $n+1$. The program is given at the end of the sequence.}

\fiche{Fractions and fractional expressions}
\begin{cours}[Rules and conditions]
For nonzero denominators:
\[\frac ac\pm\frac bc=\frac{a\pm b}{c},\quad\frac ac\pm\frac bd=\frac{ad\pm bc}{cd},\quad\frac ac\times\frac bd=\frac{ab}{cd}.\]
A fraction $c/d$ has an inverse only if $c\ne0$ and $d\ne0$. Then
\[\frac ab\div\frac cd=\frac ab\times\frac dc=\frac{ad}{bc}\quad(b,c,d\ne0).\]
\end{cours}
\ex{Calculate and reduce to lowest terms:
\[A=\left(\frac23+\frac16\right)\times\frac45,\quad B=\frac12+\frac34\times\frac56,\quad C=\frac45\div\frac23\times\frac7{10},\]
\[D=\frac{\left(\frac25+\frac13\right)\times\frac34}{\frac12-\frac16}.\]}
\reponse[5]{$A=\frac56\times\frac45=\frac23$. $B=\frac12+\frac58=\frac98$. $C=\frac45\times\frac32\times\frac7{10}=\frac{21}{25}$. $D=\frac{11}{15}\times\frac34\div\frac13=\frac{11}{20}\times3=\frac{33}{20}$. In $C$, multiplication and division are carried out from left to right.}
\ex{Simplify $\dfrac{x^2-1}{x-1}$, stating the domain. Can we then take $x=1$?}
\reponse[3.2]{For $x\ne1$, $\frac{(x-1)(x+1)}{x-1}=x+1$. The original expression remains undefined at $1$; the simplification does not lift the restriction.}

\fiche{Square roots and validity conditions}
\begin{cours}[Definition and properties]
For $a\geq0$, $\sqrt a$ is the unique positive or zero number whose square is $a$. Thus $(\sqrt a)^2=a$. For every $x\in\R$, $\sqrt{x^2}=|x|$.
For $a,b\geq0$, $\sqrt{ab}=\sqrt a\sqrt b$. If $b>0$, $\sqrt{a/b}=\sqrt a/\sqrt b$.
\end{cours}
\preuve[3.2]{The number $\sqrt a\sqrt b$ is positive or zero and its square is $ab$; by uniqueness of the positive square root, it equals $\sqrt{ab}$. The same argument applied to the quotient establishes the second formula, under $b>0$. Similarly $|x|\geq0$ and $|x|^2=x^2$, so $\sqrt{x^2}=|x|$.}
\ex{Simplify $\sqrt{18}$; $\sqrt2\sqrt8$; $\sqrt{7/25}$; $\sqrt{27/75}$; $\sqrt{(-5)^2}$.}
\reponse[3.2]{$3\sqrt2$; $4$; $\sqrt7/5$; $\sqrt{9/25}=3/5$; $5$.}
\vigilance{There is no general rule $\sqrt{a+b}=\sqrt a+\sqrt b$. With $a=4,b=9$, we would wrongly have $\sqrt{13}=5$. Saying these expressions are always different would also be wrong, for example if $a=0$.}
\ex{Simplify $2\sqrt{18}-\sqrt{50}+\sqrt8$. Bound $\sqrt{20}$ between two consecutive integers.}
\reponse[3.2]{$6\sqrt2-5\sqrt2+2\sqrt2=3\sqrt2$. Since $16<20<25$, $4<\sqrt{20}<5$.}

\fiche{First-degree equations}
\begin{cours}[Definition]
To solve an equation is to find all the values of the unknown that satisfy the equality. Two equations are equivalent if they have the same solutions. Adding the same expression or multiplying by the same nonzero number preserves equivalence.
\end{cours}
\begin{methode}
Expand if necessary, collect the terms in $x$, then isolate $x$. If its coefficient is zero, examine the remaining equality: it may be always true or impossible. Check in the original equation.
\end{methode}
\ex{Solve:
\begin{align*}
a)&\ 5x-3=2x+9,& b)&\ 3(2x-1)=2(x+4)+7,\\
c)&\ -2(x+3)=4x-(x-1),&d)&\ \frac{x+2}{3}=x-1,\\
e)&\ \frac{-5x+5}{8}=-\frac34,&f)&\ \frac{-3x-1}{9}=-\frac74.
\end{align*}}
\reponse[6]{a) $3x=12$, so $x=4$. b) $6x-3=2x+15$, so $x=9/2$. c) $-2x-6=3x+1$, so $x=-7/5$. d) $x+2=3x-3$, so $x=5/2$. e) $-5x+5=-6$, so $x=11/5$. f) $4(-3x-1)=-63$, so $-12x=-59$ and $x=59/12$.}
\ex{Solve $2(x+1)=2x+2$, then $2(x+1)=2x+3$.}
\reponse[2.4]{The first reduces to $2=2$: $S=\R$. The second gives $2=3$: $S=\varnothing$. We do not divide by a zero coefficient.}

\fiche{Zero-product property and equations with a quotient}
\begin{cours}[Properties]
$AB=0\iff A=0$ or $B=0$. For $B\ne0$, $A/B=0\iff A=0$. To solve $A(x)/B(x)=k$, first require $B(x)\ne0$, then solve $A(x)=kB(x)$ and discard the forbidden values.
\end{cours}
\ex{Solve $(4x+6)(3-7x)=0$; $3x(-7x+5)=0$; $\dfrac{2x+1}{3x-5}=0$.}
\reponse[3.2]{Solutions: $\{-3/2;3/7\}$; $\{0;5/7\}$; $\{-1/2\}$, with forbidden value $5/3$ in the last case.}
\ex{Factorise then solve $(1-6x)^2-(1-6x)(2+5x)=0$ and $(4x-1)(x+6)+(4x-1)=0$.}
\reponse[3.2]{First: $(6x-1)(11x+1)=0$, so $S=\{1/6;-1/11\}$. Second: $(4x-1)(x+7)=0$, so $S=\{1/4;-7\}$.}
\ex{Solve $\dfrac{x+1}{x-2}=3$, then $\dfrac{x^2-4}{x-2}=4$.}
\reponse[4]{In both cases, $x\ne2$. First: $x+1=3x-6$, so $x=7/2$, allowed. Second: for $x\ne2$, the quotient equals $x+2$; the equation requires $x=2$, forbidden. So no solution.}

\fiche{Equations with a square; isolating a variable}
\begin{cours}[Reminder]
In $\R$, $x^2=a$ has two solutions $\pm\sqrt a$ if $a>0$, one solution $0$ if $a=0$, and none if $a<0$.
\end{cours}
\ex{Solve $x^2=0$; $x^2=-8$; $x^2=7$; $25x^2=81$; $(x+2)^2=9$.}
\reponse[4]{In order: $\{0\}$; $\varnothing$; $\{-\sqrt7;\sqrt7\}$; $\{-9/5;9/5\}$; $\{-5;1\}$. For the last one, $x+2=-3$ or $x+2=3$.}
\begin{methode}[Isolating a variable in a formula]
Describe the operations performed on both sides; state the conditions that allow dividing or taking a square root.
\end{methode}
\ex{Isolate $h$ in $V=\pi r^2h$ with $r>0$; isolate $v$ in $E=\frac12mv^2$ with $m>0$, $E\geq0$ and $v\geq0$; isolate $x$ in $y=ax+b$ with $a\ne0$.}
\reponse[4]{$h=V/(\pi r^2)$. Then $v^2=2E/m$; since $v\geq0$, $v=\sqrt{2E/m}$. Finally $x=(y-b)/a$. Without the condition $v\geq0$, two signs would be possible for $E>0$.}
\rituel{Give the domain of $1/(x+3)$; calculate $\sqrt{49}$; the slope between $(0;2)$ and $(3;8)$.}{$\R\setminus\{-3\}$; $7$; $2$.}

\fiche{Inequalities and comparison}
\begin{cours}[Rules]
Inequalities of the same direction can be added: $a\leq b$, $c\leq d\Rightarrow a+c\leq b+d$. Adding the same real number preserves the direction. Multiplying or dividing by a strictly positive real number preserves the direction; by a strictly negative one, it reverses it.
Comparing $A$ and $B$ amounts to studying the sign of $A-B$. If $A,B>0$, one can also compare $A/B$ to $1$.
\end{cours}
\ex{1. If $x\leq-4$, what can be said about $3x-4$? 2. If $x>-6$, what can be said about $2x-3$? 3. If $7\leq x\leq16$ and $3\leq y\leq12$, bound $8x+7y$. 4. If $-10\leq x\leq-6$, bound $2x-6$.}
\reponse[4]{1. $3x-4\leq-16$. 2. $2x-3>-15$. 3. $56\leq8x\leq128$ and $21\leq7y\leq84$, so $77\leq8x+7y\leq212$. The bounds are attainable. 4. $-26\leq2x-6\leq-18$.}
\ex{Solve $3x+1\geq x-7$; $2x-11\leq4x-5$; $6x+19\leq2x-3$.}
\reponse[4]{1. $2x\geq-8$, that is $[-4;+\infty[$. 2. $-6\leq2x$, that is $[-3;+\infty[$. 3. $4x\leq-22$, that is $]-\infty;-11/2]$.}

\fiche{Sign of an affine expression and of a product}
\begin{cours}[Property]
For $a\ne0$, $ax+b$ vanishes at $x_0=-b/a$; its sign to the right is that of $a$, to the left the opposite. If $a=0$, the sign is that of the constant $b$. For a product, multiply the signs of the factors.
\end{cours}
\ex{Draw the sign tables of $3-x$, $2x-6$, then $P(x)=(2x-6)(x+1)$.}
\reponse[5]{Both affine expressions vanish at $3$: $3-x$ is positive on the left, negative on the right; $2x-6$ has the opposite signs.
\begin{center}\renewcommand{\arraystretch}{1.5}\begin{tabular}{c|ccccccc}
$x$&$-\infty$&&$-1$&&$3$&&$+\infty$\\\hline
$2x-6$&&$-$&$-$&$-$&$0$&$+$&\\
$x+1$&&$-$&$0$&$+$&$+$&$+$&\\\hline
$P(x)$&&$+$&$0$&$-$&$0$&$+$&
\end{tabular}\end{center}}
\ex{Solve $(7+3x)(2x+6)>0$.}
\reponse[5]{The zeros are $-7/3$ and $-3$, to be ordered $-3<-7/3$. The product is positive outside these roots and negative between them. The inequality is strict: $S=]-\infty;-3[\cup]-7/3;+\infty[$.}

\fiche{Sign of a quotient and algorithmic synthesis}
\begin{methode}
Find the forbidden values, then the zeros of the numerator. Order all the critical values and apply the sign rule. On the quotient row, a double bar indicates a forbidden value.
\end{methode}
\ex{Draw the sign table of $Q(x)=\dfrac{-4x+6}{x-3}$, then solve $\dfrac{2x-1}{-4x+3}\leq0$.}
\reponse[5]{For $Q$, zero $3/2$, forbidden $3$. Signs: negative on $]-\infty;3/2[$, zero at $3/2$, positive on $]3/2;3[$, undefined at $3$, negative on $]3;+\infty[$.
For the second expression, zero $1/2$, forbidden $3/4$. The signs are $-,+,-$ on the three intervals. Thus $S=]-\infty;1/2]\cup]3/4;+\infty[$.}
\begin{lstlisting}
def premiere_puissance(seuil):
    # seuil reel >= 1
    n, p = 0, 1
    while p <= seuil:
        n = n+1
        p = 2*p
    return n, p
\end{lstlisting}
\ex{What does \texttt{premiere\_puissance(100)} return? What relation between \texttt{p} and \texttt{n} stays true at each step? How can the sign of a number then be tested in Python?}
\reponse[3.2]{The result is $(7,128)$. The invariant is $p=2^n$, true at initialisation and then preserved. For the sign, first test \texttt{x > 0}, then \texttt{x < 0}, otherwise $x=0$. Do not replace the loop's non-strict inequality by a strict one without changing the problem posed.}
